% p1-chp02-lexer.tex

\chapter{Lexical Analysis and the Lexer}
\label{chp:lexer}

The lexical analyzer, or lexer, is the first phase of the compiler. It reads
the source program as a stream of characters and groups them into tokens,
the units that the parser works with \cite{aho2006compilers}. This chapter
defines the full token specification of OctaC, the deterministic finite
automata that recognize those tokens
\cite{aho2006compilers,hopcroft2006automata}, and the algorithm the lexer
uses to drive them. The source code of the lexer is available at:
\url{https://github.com/Octa-C/lexer}

\section{Token Specification}

This section lists the complete token set for the lexical analyzer, in
Table~\ref{tab:token-set}. Every token belongs to a category, and some are
divided further into a subcategory. The loop keywords \code{for}, \code{do},
\code{until}, and \code{continue}, for example, belong to the Keyword category
and the Loops subcategory. The Token column gives the token as it is written
in the source, or a descriptive name where one token covers several lexemes.
The Values column is filled only for those tokens that match more than one
lexeme, or a pattern, such as the scalar datatypes or the integer literals,
and lists the lexemes or the regular expression they match. The last column
gives the generated name of the token, built from its category, subcategory,
and token in the form \code{CATEGORY\_SUBCATEGORY\_TOKEN}, which is the name
used to identify the token from here on.

The logical and equality operators are written as English words, so they are
listed among the keywords rather than the operators. The two compiler tokens,
EOF and Error, do not correspond to any source text and are produced by the
lexer itself.

\begin{table}[tbp]
\centering
\fontsize{6.2}{7.2}\selectfont
\renewcommand{\arraystretch}{1.1}
\setlength{\tabcolsep}{2pt}
\begin{tabular}{@{}ll@{\hspace{6pt}}ll@{\hspace{6pt}}l@{}}
\toprule
\textbf{Category} & \textbf{Subcategory} & \textbf{Token} & \textbf{Values} & \textbf{Generated token} \\
\midrule
Keyword & Loops & \texttt{for} &  & \tok{KW_LOOPS_FOR} \\
Keyword & Loops & \texttt{do} &  & \tok{KW_LOOPS_DO} \\
Keyword & Loops & \texttt{until} &  & \tok{KW_LOOPS_UNTIL} \\
Keyword & Loops & \texttt{continue} &  & \tok{KW_LOOPS_CONTINUE} \\
Keyword & Function & \texttt{function} &  & \tok{KW_FUNC_FUNCTION} \\
Keyword & Function & \texttt{return} &  & \tok{KW_FUNC_RETURN} \\
Keyword & Conditional & \texttt{if} &  & \tok{KW_CONDITIONAL_IF} \\
Keyword & Conditional & \texttt{else} &  & \tok{KW_CONDITIONAL_ELSE} \\
Keyword & Conditional & \texttt{elseif} &  & \tok{KW_CONDITIONAL_ELSEIF} \\
Keyword & Conditional & \texttt{caseof} &  & \tok{KW_CONDITIONAL_CASEOF} \\
Keyword & Conditional & \texttt{case} &  & \tok{KW_CONDITIONAL_CASE} \\
Keyword & Conditional & \texttt{default} &  & \tok{KW_CONDITIONAL_DEFAULT} \\
Keyword & Word-based Logic/Rel Ops & \texttt{and} &  & \tok{KW_WORDBASED_LOGICREL_OPS_AND} \\
Keyword & Word-based Logic/Rel Ops & \texttt{or} &  & \tok{KW_WORDBASED_LOGICREL_OPS_OR} \\
Keyword & Word-based Logic/Rel Ops & \texttt{not} &  & \tok{KW_WORDBASED_LOGICREL_OPS_NOT} \\
Keyword & Word-based Logic/Rel Ops & \texttt{is} &  & \tok{KW_WORDBASED_LOGICREL_OPS_IS} \\
Keyword & Word-based Logic/Rel Ops & \texttt{in} &  & \tok{KW_WORDBASED_LOGICREL_OPS_IN} \\
Keyword & Misc & \texttt{break} &  & \tok{KW_MISC_BREAK} \\
\midrule
Datatype &  & Scalar & \texttt{i64}, \texttt{i32}, \texttt{i16}, \texttt{f64}, \texttt{f32}, \texttt{f16}, \texttt{bool} & \tok{DT_SCALAR} \\
Datatype &  & \texttt{vector} &  & \tok{DT_VECTOR} \\
Datatype &  & \texttt{matrix} &  & \tok{DT_MATRIX} \\
Datatype &  & \texttt{string} &  & \tok{DT_STRING} \\
\midrule
Punctuation &  & \texttt{\{} &  & \tok{PUNCTUATION_LBRACE} \\
Punctuation &  & \texttt{\}} &  & \tok{PUNCTUATION_RBRACE} \\
Punctuation &  & \texttt{(} &  & \tok{PUNCTUATION_LPAREN} \\
Punctuation &  & \texttt{)} &  & \tok{PUNCTUATION_RPAREN} \\
Punctuation &  & \texttt{[} &  & \tok{PUNCTUATION_LBRACKET} \\
Punctuation &  & \texttt{]} &  & \tok{PUNCTUATION_RBRACKET} \\
Punctuation &  & \texttt{;} &  & \tok{PUNCTUATION_SEMICOLON} \\
Punctuation &  & \texttt{->} &  & \tok{PUNCTUATION_ARROW} \\
Punctuation &  & \texttt{,} &  & \tok{PUNCTUATION_COMMA} \\
Punctuation &  & \texttt{...} &  & \tok{PUNCTUATION_ELLIPSIS} \\
Punctuation &  & \texttt{:} &  & \tok{PUNCTUATION_COLON} \\
\midrule
Operators & Standard & \texttt{+-} & \texttt{+}, \texttt{-} & \tok{OPERATORS_STANDARD_PLUSMINUS} \\
Operators & Standard & \texttt{* / \%} & \texttt{*}, \texttt{/}, \texttt{\%} & \tok{OPERATORS_STANDARD_STAR_SLASH_MOD} \\
Operators & Matrix & \texttt{'} &  & \tok{OPERATORS_MATRIX_TRANSPOSE} \\
Operators & Matrix & \texttt{.*} &  & \tok{OPERATORS_MATRIX_DOT_STAR} \\
Operators & Matrix & \texttt{./} &  & \tok{OPERATORS_MATRIX_DOT_SLASH} \\
Operators & Relational and Comparison Ops & \texttt{<} &  & \tok{OPERATORS_RELATIONAL_AND_COMPARISON_OPS_LESS} \\
Operators & Relational and Comparison Ops & \texttt{>} &  & \tok{OPERATORS_RELATIONAL_AND_COMPARISON_OPS_GREATER} \\
Operators & Relational and Comparison Ops & \texttt{<=} &  & \tok{OPERATORS_RELATIONAL_AND_COMPARISON_OPS_LESS_EQ} \\
Operators & Relational and Comparison Ops & \texttt{>=} &  & \tok{OPERATORS_RELATIONAL_AND_COMPARISON_OPS_GREATER_EQ} \\
Operators & Assignment &  & \texttt{=}, \texttt{+=}, \texttt{-=}, \texttt{*=}, \texttt{/=} & \tok{OPERATORS_ASSIGNMENT} \\
\midrule
Identifier &  & identifier & \texttt{[a-zA-Z][a-zA-Z0-9\_]*} & \tok{IDENTIFIER} \\
\midrule
Constants &  & bool constant & \texttt{1}, \texttt{0} & \tok{CONSTANTS_BOOL_CONSTANT} \\
Constants &  & int literal & \texttt{[0-9]+} & \tok{CONSTANTS_INT_LITERAL} \\
Constants &  & float literal & \begin{tabular}[t]{@{}l@{}}\texttt{[0-9]+\textbackslash .[0-9]+}\\ \texttt{[0-9]+(\textbackslash .[0-9]+)?[eE][+-]?[0-9]+}\end{tabular} & \tok{CONSTANTS_FLOAT_LITERAL} \\
Constants &  & string literal & \texttt{"[\textasciicircum "\textbackslash n]*"} & \tok{CONSTANTS_STRING_LITERAL} \\
\midrule
Compiler &  & EOF &  & \tok{COMPILER_EOF} \\
Compiler &  & Error &  & \tok{COMPILER_ERROR} \\
\bottomrule
\end{tabular}
\caption{The OctaC token set.}
\label{tab:token-set}
\end{table}

\section{Deterministic Finite Automata}

The lexer recognizes every token of OctaC with a single deterministic finite
automaton, driven by a transition table
\cite{aho2006compilers,hopcroft2006automata}. The automaton only recognizes
the shape of a lexeme, such as a word, a number, or an operator. Which token
that lexeme becomes is decided afterwards by a lookup, described in
Section~\ref{sec:lexer-algorithm}. This keeps the automaton small: reserved
words do not need a path of their own, because they have the same shape as
identifiers. The subsections below describe the part of the automaton that
handles each class of tokens, and each has its own diagram. The last
subsection merges them into the one automaton that the lexer runs.

\subsection{Character Classes}

The automaton does not look at individual characters. Each input byte is first
mapped to one of 17 character classes, and the transition table has one
column per class. The letters \code{e} and \code{E} have a class of their own
because they can begin the exponent of a number, while everywhere else they
behave like any other letter. Table~\ref{tab:classes} lists the classes.

\begin{table}[H]
\centering
\footnotesize
\begin{tabular}{@{}ll@{}}
\code{LETTER} & \code{A-D} \code{F-Z} \code{a-d} \code{f-z} \\
\code{EXP} & \code{E} \code{e} \\
\code{DIGIT} & \code{0-9} \\
\code{\_} & \code{\_} \\
\code{WS} & \code{\symbol{92}t} \code{\symbol{92}v} \code{\symbol{92}f} \code{\symbol{92}r} \code{space} \\
\code{NL} & \code{\symbol{92}n} \\
\code{QUOTE} & \code{\symbol{34}} \\
\code{DOT} & \code{.} \\
\code{PLUS} & \code{+} \\
\code{MINUS} & \code{-} \\
\code{STAR} & \code{*} \\
\code{SLASH} & \code{/} \\
\code{LT} & \code{<} \\
\code{GT} & \code{>} \\
\code{EQ} & \code{=} \\
\code{SIMPLE} & \code{\%} \code{\textquotesingle{}} \code{(} \code{)} \code{,} \code{:} \code{;} \code{[} \code{]} \code{\symbol{123}} \code{\symbol{125}} \\
\code{OTHER} & any other character \\
\end{tabular}
\caption{Character classes of the lexer.}
\label{tab:classes}
\end{table}

\subsection{States}

The automaton has 19 states, listed in Table~\ref{tab:dfa-states}. A state is
accepting when it has an action, which says what the lexeme read so far
produces. States without an action are only stepping stones towards a longer
lexeme.

\begin{table}[H]
\centering
\footnotesize
\renewcommand{\arraystretch}{1.1}
\begin{tabular}{@{}l>{\raggedright\arraybackslash}p{190pt}>{\raggedright\arraybackslash}p{130pt}@{}}
\toprule
\textbf{State} & \textbf{Reached after reading} & \textbf{Action} \\
\midrule
\code{DEAD} & no valid continuation exists (error state) & none \\
\code{START} & nothing yet & none \\
\code{IDENT} & \code{[a-zA-Z][a-zA-Z0-9\_]*} & keyword lookup, else identifier \\
\code{INT} & \code{[0-9]+} & integer literal \\
\code{INT\_DOT} & an integer and a dot, such as \code{12.} & none \\
\code{FLOAT} & \code{[0-9]+\textbackslash .[0-9]+} & float literal \\
\code{EXP} & a number and \code{e} or \code{E} & none \\
\code{EXP\_SIGN} & a number, \code{e} or \code{E}, and a sign & none \\
\code{FLOAT\_EXP} & a number with a complete exponent & float literal \\
\code{OP\_DONE} & an operator or punctuator that cannot be extended, such as \code{\{}, \code{->}, \code{<=}, \code{...}, \code{.*} or \code{=} & operator lookup \\
\code{MINUS} & \code{-}, which may extend to \code{->} or \code{-=} & operator lookup \\
\code{OP\_EQ} & \code{<}, \code{>}, \code{+} or \code{*}, which may extend with \code{=} & operator lookup \\
\code{SLASH} & \code{/}, which may extend to \code{/=} or begin a comment & operator lookup \\
\code{COMMENT} & \code{//} and the rest of the line & skip \\
\code{DOT} & \code{.}, waiting for \code{*}, \code{/} or another \code{.} & none \\
\code{DOT\_DOT} & \code{..}, waiting for a third \code{.} & none \\
\code{STR\_BODY} & an opening quote and the text after it & none \\
\code{STR\_END} & a complete string literal & string literal \\
\code{WS} & a run of whitespace & skip \\
\bottomrule
\end{tabular}
\caption{States of the lexer DFA.}
\label{tab:dfa-states}
\end{table}

\subsection{Identifiers, Keywords, and Type Names}

An identifier has the pattern \code{[a-zA-Z][a-zA-Z0-9\_]*}: a letter followed
by any number of letters, digits and underscores. The DFA has one accepting
state, \code{IDENT}, and it emits \textless IDENTIFIER, value\textgreater.
Keywords and type names have the same shape, so the DFA reads them as
identifiers as well. After the DFA accepts, the lexer looks the lexeme up in
Table~\ref{tab:lookup}. A lexeme that is in the table gets the keyword or type
token given there in place of \textless IDENTIFIER, value\textgreater.

An underscore cannot start an identifier, so it has no transition out of
the start state and is reported as a lexical error.

\begin{figure}[H]
\centering
\begin{adjustbox}{max width=\linewidth,max height=0.3\textheight}
\begin{tikzpicture}[x=1in,y=1in,line width=0.4pt]
\node[draw,circle,fill=white,inner sep=1pt,minimum size=0.42in] (START) at (0.9,0.7) {START};
\node[draw,circle,fill=white,inner sep=1pt,minimum size=0.42in,double,double distance=1.3pt] (IDENT) at (3.0,0.7) {IDENT};
\draw[-{Triangle[length=4bp,width=3bp]}] (0.4,0.7) -- (START);
\node[inner sep=0pt,align=left,anchor=west] (IDENT-out) at (3.6,0.7) {\textless{}IDENTIFIER, value\textgreater{}};
\draw[-{Triangle[length=4bp,width=3bp]}] (START) to[bend left=0] node[pos=0.5,above,inner sep=1pt,align=center] {\code{a-z} \code{A-Z}} (IDENT);
\path[-{Triangle[length=4bp,width=3bp]}] (IDENT) edge[loop above,min distance=7mm] node[above,inner sep=1pt,align=center] {\code{a-z} \code{A-Z} \code{0-9} \code{\_}} (IDENT);
\draw[-{Triangle[length=4bp,width=3bp]}] (IDENT) to[bend left=0] (IDENT-out.west);
\end{tikzpicture}
\end{adjustbox}
\caption{DFA for identifiers, keywords, and type names.}
\label{fig:dfa-ident}
\end{figure}

\begin{table}[H]
\centering
\footnotesize
\begin{tabular}{@{}ll@{}}
\code{for} & \textless{}KW\_LOOPS\_FOR,\textgreater{} \\
\code{do} & \textless{}KW\_LOOPS\_DO,\textgreater{} \\
\code{until} & \textless{}KW\_LOOPS\_UNTIL,\textgreater{} \\
\code{continue} & \textless{}KW\_LOOPS\_CONTINUE,\textgreater{} \\
\code{function} & \textless{}KW\_FUNC\_FUNCTION,\textgreater{} \\
\code{return} & \textless{}KW\_FUNC\_RETURN,\textgreater{} \\
\code{if} & \textless{}KW\_CONDITIONAL\_IF,\textgreater{} \\
\code{else} & \textless{}KW\_CONDITIONAL\_ELSE,\textgreater{} \\
\code{elseif} & \textless{}KW\_CONDITIONAL\_ELSEIF,\textgreater{} \\
\code{caseof} & \textless{}KW\_CONDITIONAL\_CASEOF,\textgreater{} \\
\code{case} & \textless{}KW\_CONDITIONAL\_CASE,\textgreater{} \\
\code{default} & \textless{}KW\_CONDITIONAL\_DEFAULT,\textgreater{} \\
\code{and} & \textless{}KW\_WORDBASED\_LOGICREL\_OPS\_AND,\textgreater{} \\
\code{or} & \textless{}KW\_WORDBASED\_LOGICREL\_OPS\_OR,\textgreater{} \\
\code{not} & \textless{}KW\_WORDBASED\_LOGICREL\_OPS\_NOT,\textgreater{} \\
\code{is} & \textless{}KW\_WORDBASED\_LOGICREL\_OPS\_IS,\textgreater{} \\
\code{in} & \textless{}KW\_WORDBASED\_LOGICREL\_OPS\_IN,\textgreater{} \\
\code{break} & \textless{}KW\_MISC\_BREAK,\textgreater{} \\
\code{i64} \code{i32} \code{i16} \code{f64} \code{f32} \code{f16} \code{bool} & \textless{}DT\_SCALAR, value\textgreater{} \\
\code{vector} & \textless{}DT\_VECTOR,\textgreater{} \\
\code{matrix} & \textless{}DT\_MATRIX,\textgreater{} \\
\code{string} & \textless{}DT\_STRING,\textgreater{} \\
\end{tabular}
\caption{Keywords and type names, and the tokens the lookup after IDENT gives them.}
\label{tab:lookup}
\end{table}

\subsection{Punctuation}

A punctuation token is one of the characters \code{\symbol{123}}
\code{\symbol{125}} \code{(} \code{)} \code{[} \code{]} \code{;} \code{,}
\code{:}, or \code{->}, or \code{...}. A single character goes from
\code{START} straight to \code{OP\_DONE}. \code{->} goes through \code{MINUS},
and \code{...} goes through \code{DOT} and \code{DOT\_DOT}. These three states
are not accepting in this automaton, because \code{-} and \code{.} on their
own are not punctuation. \code{OP\_DONE} is accepting, and
Table~\ref{tab:punct} gives the token for each lexeme.

\begin{table}[H]
\centering
\footnotesize
\begin{tabular}{@{}ll@{}}
\code{\symbol{123}} & \textless{}PUNCTUATION\_LBRACE,\textgreater{} \\
\code{\symbol{125}} & \textless{}PUNCTUATION\_RBRACE,\textgreater{} \\
\code{(} & \textless{}PUNCTUATION\_LPAREN,\textgreater{} \\
\code{)} & \textless{}PUNCTUATION\_RPAREN,\textgreater{} \\
\code{[} & \textless{}PUNCTUATION\_LBRACKET,\textgreater{} \\
\code{]} & \textless{}PUNCTUATION\_RBRACKET,\textgreater{} \\
\code{;} & \textless{}PUNCTUATION\_SEMICOLON,\textgreater{} \\
\code{->} & \textless{}PUNCTUATION\_ARROW,\textgreater{} \\
\code{,} & \textless{}PUNCTUATION\_COMMA,\textgreater{} \\
\code{...} & \textless{}PUNCTUATION\_ELLIPSIS,\textgreater{} \\
\code{:} & \textless{}PUNCTUATION\_COLON,\textgreater{} \\
\end{tabular}
\caption{Punctuation lexemes and the tokens they give.}
\label{tab:punct}
\end{table}

\begin{figure}[H]
\centering
\begin{adjustbox}{max width=\linewidth,max height=0.4\textheight}
\begin{tikzpicture}[x=1in,y=1in,line width=0.4pt]
\node[draw,circle,fill=white,inner sep=1pt,minimum size=0.42in] (START) at (0.8,2.4) {START};
\node[draw,circle,fill=white,inner sep=1pt,minimum size=0.42in,double,double distance=1.3pt] (OP_DONE) at (4.3,2.4) {OP\_DONE};
\node[draw,circle,fill=white,inner sep=1pt,minimum size=0.42in] (MINUS) at (2.3,0.9) {MINUS};
\node[draw,circle,fill=white,inner sep=1pt,minimum size=0.42in] (DOT) at (2.1,3.9) {DOT};
\node[draw,circle,fill=white,inner sep=1pt,minimum size=0.42in] (DOT_DOT) at (3.4,3.9) {DOT\_DOT};
\draw[-{Triangle[length=4bp,width=3bp]}] (0.30000000000000004,2.4) -- (START);
\node[inner sep=0pt,align=left,anchor=west] (OP_DONE-out) at (5.0,2.4) {Table~\ref{tab:punct}};
\draw[-{Triangle[length=4bp,width=3bp]}] (START) to[bend left=0] node[pos=0.5,above,inner sep=1pt,align=center] {\code{(} \code{)} \code{,} \code{:} \code{;} \code{[} \code{]} \code{\symbol{123}} \code{\symbol{125}}} (OP_DONE);
\draw[-{Triangle[length=4bp,width=3bp]}] (START) to[bend left=0] node[pos=0.5,below left,inner sep=1pt,align=center] {\code{-}} (MINUS);
\draw[-{Triangle[length=4bp,width=3bp]}] (START) to[bend left=0] node[pos=0.5,above left,inner sep=1pt,align=center] {\code{.}} (DOT);
\draw[-{Triangle[length=4bp,width=3bp]}] (MINUS) to[bend left=0] node[pos=0.5,above left,inner sep=1pt,align=center] {\code{>}} (OP_DONE);
\draw[-{Triangle[length=4bp,width=3bp]}] (DOT) to[bend left=0] node[pos=0.5,above,inner sep=1pt,align=center] {\code{.}} (DOT_DOT);
\draw[-{Triangle[length=4bp,width=3bp]}] (DOT_DOT) to[bend left=0] node[pos=0.5,right,inner sep=1pt,align=center] {\code{.}} (OP_DONE);
\draw[-{Triangle[length=4bp,width=3bp]}] (OP_DONE) to[bend left=0] (OP_DONE-out.west);
\end{tikzpicture}
\end{adjustbox}
\caption{DFA for punctuation.}
\label{fig:dfa-punct}
\end{figure}

\subsection{Operators}
\label{sec:ops}

The operators are \code{+} \code{-} \code{*} \code{/} \code{\%}
\code{\textquotesingle{}}, the matrix operators \code{.*} and \code{./}, the
comparisons \code{<} \code{>} \code{<=} \code{>=}, and the assignments
\code{=} \code{+=} \code{-=} \code{*=} \code{/=}. \code{\%}
\code{\textquotesingle{}} \code{=} go from \code{START} straight to
\code{OP\_DONE}. \code{+} \code{*} \code{<} \code{>} go to \code{OP\_EQ},
\code{-} to \code{MINUS} and \code{/} to \code{SLASH}. These three states are
accepting, because each of those characters is an operator on its own, and
they continue to \code{OP\_DONE} on \code{=}. \code{DOT} is not accepting,
because \code{.} alone is not an operator. It continues to \code{OP\_DONE} on
\code{*} or \code{/}. The state a lexeme ends in does not fix its token, so
Table~\ref{tab:ops} gives the token for each lexeme.

There is no symbolic operator for equality or inequality, since these are
written with the words \code{is} and \code{is not}.

\begin{table}[H]
\centering
\footnotesize
\begin{tabular}{@{}ll@{}}
\code{+} \code{-} & \textless{}OPERATORS\_STANDARD\_PLUSMINUS, value\textgreater{} \\
\code{*} \code{/} \code{\%} & \textless{}OPERATORS\_STANDARD\_STAR\_SLASH\_MOD, value\textgreater{} \\
\code{\textquotesingle{}} & \textless{}OPERATORS\_MATRIX\_TRANSPOSE,\textgreater{} \\
\code{.*} & \textless{}OPERATORS\_MATRIX\_DOT\_STAR,\textgreater{} \\
\code{./} & \textless{}OPERATORS\_MATRIX\_DOT\_SLASH,\textgreater{} \\
\code{<} & \textless{}OPERATORS\_RELATIONAL\_AND\_COMPARISON\_OPS\_LESS,\textgreater{} \\
\code{>} & \textless{}OPERATORS\_RELATIONAL\_AND\_COMPARISON\_OPS\_GREATER,\textgreater{} \\
\code{<=} & \textless{}OPERATORS\_RELATIONAL\_AND\_COMPARISON\_OPS\_LESS\_EQ,\textgreater{} \\
\code{>=} & \textless{}OPERATORS\_RELATIONAL\_AND\_COMPARISON\_OPS\_GREATER\_EQ,\textgreater{} \\
\code{=} \code{+=} \code{-=} \code{*=} \code{/=} & \textless{}OPERATORS\_ASSIGNMENT, value\textgreater{} \\
\end{tabular}
\caption{Operator lexemes and the tokens they give.}
\label{tab:ops}
\end{table}

\begin{figure}[H]
\centering
\begin{adjustbox}{max width=\linewidth,max height=0.6\textheight}
\begin{tikzpicture}[x=1in,y=1in,line width=0.4pt]
\node[draw,circle,fill=white,inner sep=1pt,minimum size=0.42in] (START) at (0.5,4.55) {START};
\node[draw,circle,fill=white,inner sep=1pt,minimum size=0.42in,double,double distance=1.3pt] (OP_DONE) at (5.2,4.55) {OP\_DONE};
\node[draw,circle,fill=white,inner sep=1pt,minimum size=0.42in,double,double distance=1.3pt] (MINUS) at (3.1,3.6) {MINUS};
\node[draw,circle,fill=white,inner sep=1pt,minimum size=0.42in,double,double distance=1.3pt] (OP_EQ) at (3.1,6.2) {OP\_EQ};
\node[draw,circle,fill=white,inner sep=1pt,minimum size=0.42in,double,double distance=1.3pt] (SLASH) at (3.1,2.9) {SLASH};
\node[draw,circle,fill=white,inner sep=1pt,minimum size=0.42in] (DOT) at (3.1,5.5) {DOT};
\draw[-{Triangle[length=4bp,width=3bp]}] (0.0,4.55) -- (START);
\node[inner sep=0pt,align=left,anchor=west] (OP_DONE-out) at (5.9,4.55) {Table~\ref{tab:ops}};
\node[inner sep=0pt,align=left,anchor=south] (MINUS-out) at (3.1,4.05) {Table~\ref{tab:ops}};
\node[inner sep=0pt,align=left,anchor=west] (OP_EQ-out) at (3.65,6.2) {Table~\ref{tab:ops}};
\node[inner sep=0pt,align=left,anchor=west] (SLASH-out) at (3.65,2.9) {Table~\ref{tab:ops}};
\draw[-{Triangle[length=4bp,width=3bp]}] (START) to[bend left=0] node[pos=0.5,above,inner sep=1pt,align=center] {\code{\%} \code{\textquotesingle{}} \code{=}} (OP_DONE);
\draw[-{Triangle[length=4bp,width=3bp]}] (START) to[bend left=0] node[pos=0.5,below left,inner sep=1pt,align=center] {\code{-}} (MINUS);
\draw[-{Triangle[length=4bp,width=3bp]}] (START) to[bend left=0] node[pos=0.5,above left,inner sep=1pt,align=center] {\code{*} \code{+} \code{<} \code{>}} (OP_EQ);
\draw[-{Triangle[length=4bp,width=3bp]}] (START) to[bend left=0] node[pos=0.5,below left,inner sep=1pt,align=center] {\code{/}} (SLASH);
\draw[-{Triangle[length=4bp,width=3bp]}] (START) to[bend left=0] node[pos=0.5,above left,inner sep=1pt,align=center] {\code{.}} (DOT);
\draw[-{Triangle[length=4bp,width=3bp]}] (MINUS) to[bend left=0] node[pos=0.5,above left,inner sep=1pt,align=center] {\code{=}} (OP_DONE);
\draw[-{Triangle[length=4bp,width=3bp]}] (OP_EQ) to[bend left=0] node[pos=0.5,below left,inner sep=1pt,align=center] {\code{=}} (OP_DONE);
\draw[-{Triangle[length=4bp,width=3bp]}] (SLASH) to[bend left=0] node[pos=0.5,above left,inner sep=1pt,align=center] {\code{=}} (OP_DONE);
\draw[-{Triangle[length=4bp,width=3bp]}] (DOT) to[bend left=0] node[pos=0.5,below left,inner sep=1pt,align=center] {\code{*} \code{/}} (OP_DONE);
\draw[-{Triangle[length=4bp,width=3bp]}] (OP_DONE) to[bend left=0] (OP_DONE-out.west);
\draw[-{Triangle[length=4bp,width=3bp]}] (MINUS) to[bend left=0] (MINUS-out.south);
\draw[-{Triangle[length=4bp,width=3bp]}] (OP_EQ) to[bend left=0] (OP_EQ-out.west);
\draw[-{Triangle[length=4bp,width=3bp]}] (SLASH) to[bend left=0] (SLASH-out.west);
\end{tikzpicture}
\end{adjustbox}
\caption{DFA for operators.}
\label{fig:dfa-ops}
\end{figure}

\subsection{Constants}

A constant is a bool constant, an int literal, a float literal or a string
literal.

\begin{itemize}
    \item bool constants: \code{0} \code{1}
    \item int literal: \code{[0-9]+}
    \item float literal: \code{[0-9]+\symbol{92}.[0-9]+} or
    \code{[0-9]+(\symbol{92}.[0-9]+)?[eE][+-]?[0-9]+}
    \item string literal: \code{\symbol{34}[\symbol{94}\symbol{34}\symbol{92}n]*\symbol{34}}
\end{itemize}

Digits are read in \code{INT}. A dot followed by digits leads through
\code{INT\_DOT} to \code{FLOAT}. An \code{e} or \code{E} after \code{INT} or
\code{FLOAT} leads to \code{EXP}, an optional sign to \code{EXP\_SIGN}, and
digits to \code{FLOAT\_EXP}. A quote starts \code{STR\_BODY}, which loops on
every character except a newline or a quote, and the closing quote reaches
\code{STR\_END}. \code{INT}, \code{FLOAT}, \code{FLOAT\_EXP} and
\code{STR\_END} are accepting. The lexer emits \code{0} and \code{1} as int
literals, so the bool constant is never emitted.

\code{INT\_DOT}, \code{EXP}, and \code{EXP\_SIGN} are not accepting, so a
partial literal such as \code{1.} or \code{1e+} is not a number, and the lexer
falls back to the longest number that was complete, as described in
Section~\ref{sec:lexer-algorithm}. A float must have digits on both sides of
the dot, so \code{1.} and \code{.5} are lexical errors. A digit followed by a
letter other than \code{e} or \code{E} has no transition, so \code{12abc} is
read as the integer \code{12} followed by the identifier \code{abc}. There
are no escape sequences in strings, and a string cannot span lines: a newline
inside a string has no transition, so a string that is missing its closing
quote never reaches an accepting state and is reported as an error.

\begin{figure}[H]
\centering
\begin{adjustbox}{max width=\linewidth,max height=0.42\textheight}
\begin{tikzpicture}[x=1in,y=1in,line width=0.4pt]
\node[draw,circle,fill=white,inner sep=1pt,minimum size=0.42in] (START) at (0.6,3.1) {START};
\node[draw,circle,fill=white,inner sep=1pt,minimum size=0.42in,double,double distance=1.3pt] (INT) at (1.9,4.5) {INT};
\node[draw,circle,fill=white,inner sep=1pt,minimum size=0.42in] (INT_DOT) at (3.5,4.5) {INT\_DOT};
\node[draw,circle,fill=white,inner sep=1pt,minimum size=0.42in,double,double distance=1.3pt] (FLOAT) at (5.0,4.5) {FLOAT};
\node[draw,circle,fill=white,inner sep=1pt,minimum size=0.42in] (EXP) at (4.3,3.1) {EXP};
\node[draw,circle,fill=white,inner sep=1pt,minimum size=0.42in] (EXP_SIGN) at (3.1,1.9) {EXP\_SIGN};
\node[draw,circle,fill=white,inner sep=1pt,minimum size=0.42in,double,double distance=1.3pt] (FLOAT_EXP) at (5.0,1.9) {FLOAT\_EXP};
\node[draw,circle,fill=white,inner sep=1pt,minimum size=0.42in] (STR_BODY) at (1.9,0.9) {STR\_BODY};
\node[draw,circle,fill=white,inner sep=1pt,minimum size=0.42in,double,double distance=1.3pt] (STR_END) at (4.1,0.9) {STR\_END};
\draw[-{Triangle[length=4bp,width=3bp]}] (0.09999999999999998,3.1) -- (START);
\node[inner sep=0pt,align=left,anchor=south] (INT-out) at (1.9,5.05) {\textless{}CONSTANTS\_INT\_LITERAL, value\textgreater{}};
\node[inner sep=0pt,align=left,anchor=south] (FLOAT-out) at (5.0,5.05) {\textless{}CONSTANTS\_FLOAT\_LITERAL, value\textgreater{}};
\node[inner sep=0pt,align=left,anchor=north west] (FLOAT_EXP-out) at (5.2,1.3) {\textless{}CONSTANTS\_FLOAT\_LITERAL, value\textgreater{}};
\node[inner sep=0pt,align=left,anchor=north west] (STR_END-out) at (4.25,0.4) {\textless{}CONSTANTS\_STRING\_LITERAL, value\textgreater{}};
\draw[-{Triangle[length=4bp,width=3bp]}] (START) to[bend left=0] node[pos=0.5,above left,inner sep=1pt,align=center] {\code{0-9}} (INT);
\draw[-{Triangle[length=4bp,width=3bp]}] (START) to[bend left=0] node[pos=0.5,right,inner sep=1pt,align=center] {\code{\symbol{34}}} (STR_BODY);
\path[-{Triangle[length=4bp,width=3bp]}] (INT) edge[loop left,min distance=7mm] node[left,inner sep=1pt,align=center] {\code{0-9}} (INT);
\draw[-{Triangle[length=4bp,width=3bp]}] (INT) to[bend left=0] node[pos=0.5,above,inner sep=1pt,align=center] {\code{.}} (INT_DOT);
\draw[-{Triangle[length=4bp,width=3bp]}] (INT) to[bend left=0] node[pos=0.55,below left,inner sep=1pt,align=center] {\code{E} \code{e}} (EXP);
\draw[-{Triangle[length=4bp,width=3bp]}] (INT_DOT) to[bend left=0] node[pos=0.5,above,inner sep=1pt,align=center] {\code{0-9}} (FLOAT);
\path[-{Triangle[length=4bp,width=3bp]}] (FLOAT) edge[loop right,min distance=7mm] node[right,inner sep=1pt,align=center] {\code{0-9}} (FLOAT);
\draw[-{Triangle[length=4bp,width=3bp]}] (FLOAT) to[bend left=0] node[pos=0.5,right,inner sep=1pt,align=center] {\code{E} \code{e}} (EXP);
\draw[-{Triangle[length=4bp,width=3bp]}] (EXP) to[bend left=0] node[pos=0.5,below right,inner sep=1pt,align=center] {\code{+} \code{-}} (EXP_SIGN);
\draw[-{Triangle[length=4bp,width=3bp]}] (EXP) to[bend left=0] node[pos=0.5,right,inner sep=1pt,align=center] {\code{0-9}} (FLOAT_EXP);
\draw[-{Triangle[length=4bp,width=3bp]}] (EXP_SIGN) to[bend left=0] node[pos=0.5,above,inner sep=1pt,align=center] {\code{0-9}} (FLOAT_EXP);
\path[-{Triangle[length=4bp,width=3bp]}] (FLOAT_EXP) edge[loop right,min distance=7mm] node[right,inner sep=1pt,align=center] {\code{0-9}} (FLOAT_EXP);
\path[-{Triangle[length=4bp,width=3bp]}] (STR_BODY) edge[loop below,min distance=7mm] node[below,inner sep=1pt,align=center] {\code{any except \symbol{92}n \symbol{34}}} (STR_BODY);
\draw[-{Triangle[length=4bp,width=3bp]}] (STR_BODY) to[bend left=0] node[pos=0.5,above,inner sep=1pt,align=center] {\code{\symbol{34}}} (STR_END);
\draw[-{Triangle[length=4bp,width=3bp]}] (INT) to[bend left=0] (INT-out.south);
\draw[-{Triangle[length=4bp,width=3bp]}] (FLOAT) to[bend left=0] (FLOAT-out.south);
\draw[-{Triangle[length=4bp,width=3bp]}] (FLOAT_EXP) to[bend left=0] (FLOAT_EXP-out.north west);
\draw[-{Triangle[length=4bp,width=3bp]}] (STR_END) to[bend left=0] (STR_END-out.north west);
\end{tikzpicture}
\end{adjustbox}
\caption{DFA for constants.}
\label{fig:dfa-const}
\end{figure}

\subsection{Compiler Tokens}

The end of the input emits \textless COMPILER\_EOF,\textgreater. The scanner
checks for the end of the input before it runs the DFA, so the code has no
state for it, and the transition goes straight from the start state to the
token. When the next character has no transition and no accepting state was passed, the DFA is in
\code{DEAD}, the error state of the code, and the text read so far is emitted
as \textless COMPILER\_ERROR,\textgreater.

\begin{figure}[H]
\centering
\begin{adjustbox}{max width=\linewidth,max height=0.25\textheight}
\begin{tikzpicture}[x=1in,y=1in,line width=0.4pt]
\node[draw,circle,fill=white,inner sep=1pt,minimum size=0.42in] (START) at (0.9,1.8) {START};
\node[draw,circle,fill=white,inner sep=1pt,minimum size=0.42in] (DEAD) at (3.0,0.6) {DEAD};
\draw[-{Triangle[length=4bp,width=3bp]}] (0.4,1.8) -- (START);
\node[inner sep=0pt,align=left,anchor=west] (eof-out) at (3.4,1.8) {\textless{}COMPILER\_EOF,\textgreater{}};
\node[inner sep=0pt,align=left,anchor=west] (DEAD-out) at (3.4,0.6) {\textless{}COMPILER\_ERROR,\textgreater{}};
\draw[-{Triangle[length=4bp,width=3bp]}] (START) to[bend left=0] node[pos=0.5,above,inner sep=1pt,align=center] {\code{end of input}} (eof-out.west);
\draw[-{Triangle[length=4bp,width=3bp]}] (DEAD) to[bend left=0] (DEAD-out.west);
\end{tikzpicture}
\end{adjustbox}
\caption{DFA for the compiler tokens.}
\label{fig:dfa-compiler}
\end{figure}

\subsection{Skipped Input}

Skipped input is whitespace and comments. Whitespace is a run of the \code{WS}
and \code{NL} classes and stays in \code{WS}. A comment starts with \code{//}
and runs up to, but not including, the next newline, staying in
\code{COMMENT}. \code{WS} and \code{COMMENT} are accepting but emit no token.
\code{SLASH} is not accepting in this automaton, because a single \code{/} is
not skipped input: it is the division operator of Section~\ref{sec:ops}.

\begin{figure}[H]
\centering
\begin{adjustbox}{max width=\linewidth,max height=0.35\textheight}
\begin{tikzpicture}[x=1in,y=1in,line width=0.4pt]
\node[draw,circle,fill=white,inner sep=1pt,minimum size=0.42in] (START) at (0.9,1.6) {START};
\node[draw,circle,fill=white,inner sep=1pt,minimum size=0.42in] (SLASH) at (2.2,0.7) {SLASH};
\node[draw,circle,fill=white,inner sep=1pt,minimum size=0.42in,double,double distance=1.3pt] (COMMENT) at (4.2,0.7) {COMMENT};
\node[draw,circle,fill=white,inner sep=1pt,minimum size=0.42in,double,double distance=1.3pt] (WS) at (3.0,2.5) {WS};
\draw[-{Triangle[length=4bp,width=3bp]}] (0.4,1.6) -- (START);
\node[inner sep=0pt,align=left,anchor=west] (COMMENT-out) at (4.85,0.7) {no token};
\node[inner sep=0pt,align=left,anchor=west] (WS-out) at (3.6,2.5) {no token};
\draw[-{Triangle[length=4bp,width=3bp]}] (START) to[bend left=0] node[pos=0.5,below left,inner sep=1pt,align=center] {\code{/}} (SLASH);
\draw[-{Triangle[length=4bp,width=3bp]}] (START) to[bend left=0] node[pos=0.5,above left,inner sep=1pt,align=center] {\code{WS} \code{NL}} (WS);
\draw[-{Triangle[length=4bp,width=3bp]}] (SLASH) to[bend left=0] node[pos=0.5,above,inner sep=1pt,align=center] {\code{/}} (COMMENT);
\path[-{Triangle[length=4bp,width=3bp]}] (COMMENT) edge[loop above,min distance=7mm] node[above,inner sep=1pt,align=center] {\code{any except \symbol{92}n}} (COMMENT);
\path[-{Triangle[length=4bp,width=3bp]}] (WS) edge[loop above,min distance=7mm] node[above,inner sep=1pt,align=center] {\code{WS} \code{NL}} (WS);
\draw[-{Triangle[length=4bp,width=3bp]}] (COMMENT) to[bend left=0] (COMMENT-out.west);
\draw[-{Triangle[length=4bp,width=3bp]}] (WS) to[bend left=0] (WS-out.west);
\end{tikzpicture}
\end{adjustbox}
\caption{DFA for whitespace and comments.}
\label{fig:dfa-skip}
\end{figure}

\subsection{The Merged DFA}

The classes are merged into one DFA, the one the lexer runs. Its states carry the
names the source gives them.

\begin{enumerate}
    \item A new start state has a transition for the first character of
    every class.
    \item Classes that share a prefix share states. Keywords, type names
    and identifiers are all read in IDENT, and operators and punctuation
    end in OP\_DONE, MINUS, OP\_EQ, SLASH, DOT and DOT\_DOT.
    \item A state that is accepting for one class stays accepting in the
    merged DFA, so MINUS and SLASH are accepting here.
    \item IDENT emits \textless IDENTIFIER, value\textgreater, which
    Table~\ref{tab:lookup} replaces when the lexeme is a keyword or type
    name. OP\_EQ and OP\_DONE emit different tokens depending on the
    lexeme. The arrow from each of them names the tables that list them.
    \item The DFA runs until the next character has no transition. The
    token of the last accepting state passed is emitted, and if none was
    passed the consumed text is a \textless COMPILER\_ERROR,\textgreater{}
    token. DEAD is not drawn: every missing transition leads to it.
\end{enumerate}

\begin{figure}[H]
\centering
\begin{adjustbox}{max width=\linewidth,max height=0.85\textheight}
\begin{tikzpicture}[x=1in,y=1in,line width=0.4pt]
\node[draw,circle,fill=white,inner sep=1pt,minimum size=0.42in] (START) at (0.6,7.1) {START};
\node[draw,circle,fill=white,inner sep=1pt,minimum size=0.42in,double,double distance=1.3pt] (IDENT) at (2.5,10.4) {IDENT};
\node[draw,circle,fill=white,inner sep=1pt,minimum size=0.42in,double,double distance=1.3pt] (INT) at (2.5,2.8) {INT};
\node[draw,circle,fill=white,inner sep=1pt,minimum size=0.42in] (INT_DOT) at (4.0,2.8) {INT\_DOT};
\node[draw,circle,fill=white,inner sep=1pt,minimum size=0.42in,double,double distance=1.3pt] (FLOAT) at (5.5,2.8) {FLOAT};
\node[draw,circle,fill=white,inner sep=1pt,minimum size=0.42in] (EXP) at (4.9,2.0) {EXP};
\node[draw,circle,fill=white,inner sep=1pt,minimum size=0.42in] (EXP_SIGN) at (3.6,1.2000000000000002) {EXP\_SIGN};
\node[draw,circle,fill=white,inner sep=1pt,minimum size=0.42in,double,double distance=1.3pt] (FLOAT_EXP) at (5.4,1.2000000000000002) {FLOAT\_EXP};
\node[draw,circle,fill=white,inner sep=1pt,minimum size=0.42in,double,double distance=1.3pt] (OP_DONE) at (4.9,7.1) {OP\_DONE};
\node[draw,circle,fill=white,inner sep=1pt,minimum size=0.42in,double,double distance=1.3pt] (MINUS) at (2.5,6.1) {MINUS};
\node[draw,circle,fill=white,inner sep=1pt,minimum size=0.42in,double,double distance=1.3pt] (OP_EQ) at (2.5,8.0) {OP\_EQ};
\node[draw,circle,fill=white,inner sep=1pt,minimum size=0.42in,double,double distance=1.3pt] (SLASH) at (2.5,5.1) {SLASH};
\node[draw,circle,fill=white,inner sep=1pt,minimum size=0.42in,double,double distance=1.3pt] (COMMENT) at (2.5,4.1) {COMMENT};
\node[draw,circle,fill=white,inner sep=1pt,minimum size=0.42in] (DOT) at (2.5,9.100000000000001) {DOT};
\node[draw,circle,fill=white,inner sep=1pt,minimum size=0.42in] (DOT_DOT) at (4.0,9.55) {DOT\_DOT};
\node[draw,circle,fill=white,inner sep=1pt,minimum size=0.42in] (STR_BODY) at (2.5,11.700000000000001) {STR\_BODY};
\node[draw,circle,fill=white,inner sep=1pt,minimum size=0.42in,double,double distance=1.3pt] (STR_END) at (4.6,11.700000000000001) {STR\_END};
\node[draw,circle,fill=white,inner sep=1pt,minimum size=0.42in,double,double distance=1.3pt] (WS) at (2.5,13.0) {WS};
\draw[-{Triangle[length=4bp,width=3bp]}] (0.09999999999999998,7.1) -- (START);
\node[inner sep=0pt,align=left,anchor=west] (IDENT-out) at (3.15,10.4) {\textless{}IDENTIFIER, value\textgreater{}};
\node[inner sep=0pt,align=left,anchor=north] (INT-out) at (1.6,2.3) {\textless{}CONSTANTS\_INT\_LITERAL, value\textgreater{}};
\node[inner sep=0pt,align=left,anchor=south] (FLOAT-out) at (5.5,3.3499999999999996) {\textless{}CONSTANTS\_FLOAT\_LITERAL, value\textgreater{}};
\node[inner sep=0pt,align=left,anchor=north] (FLOAT_EXP-out) at (5.4,0.65) {\textless{}CONSTANTS\_FLOAT\_LITERAL, value\textgreater{}};
\node[inner sep=0pt,align=left,anchor=west] (OP_DONE-out) at (5.35,7.1) {Tables~\ref{tab:punct} and~\ref{tab:ops}};
\node[inner sep=0pt,align=left,anchor=north] (MINUS-out) at (2.5,5.7) {\textless{}OPERATORS\_STANDARD\_PLUSMINUS, value\textgreater{}};
\node[inner sep=0pt,align=left,anchor=south] (OP_EQ-out) at (2.5,8.35) {Table~\ref{tab:ops}};
\node[inner sep=0pt,align=left,anchor=west] (SLASH-out) at (2.9,4.75) {\textless{}OPERATORS\_STANDARD\_STAR\_SLASH\_MOD, value\textgreater{}};
\node[inner sep=0pt,align=left,anchor=west] (COMMENT-out) at (3.1,4.1) {no token};
\node[inner sep=0pt,align=left,anchor=north] (STR_END-out) at (4.6,11.25) {\textless{}CONSTANTS\_STRING\_LITERAL, value\textgreater{}};
\node[inner sep=0pt,align=left,anchor=west] (WS-out) at (3.1,13.0) {no token};
\draw[-{Triangle[length=4bp,width=3bp]}] (IDENT) to[bend left=0] (IDENT-out.west);
\draw[-{Triangle[length=4bp,width=3bp]}] (INT) to[bend left=0] (INT-out.north);
\draw[-{Triangle[length=4bp,width=3bp]}] (FLOAT) to[bend left=0] (FLOAT-out.south);
\draw[-{Triangle[length=4bp,width=3bp]}] (FLOAT_EXP) to[bend left=0] (FLOAT_EXP-out.north);
\draw[-{Triangle[length=4bp,width=3bp]}] (OP_DONE) to[bend left=0] (OP_DONE-out.west);
\draw[-{Triangle[length=4bp,width=3bp]}] (MINUS) to[bend left=0] (MINUS-out.north);
\draw[-{Triangle[length=4bp,width=3bp]}] (OP_EQ) to[bend left=0] (OP_EQ-out.south);
\draw[-{Triangle[length=4bp,width=3bp]}] (SLASH) to[bend left=0] (SLASH-out.west);
\draw[-{Triangle[length=4bp,width=3bp]}] (COMMENT) to[bend left=0] (COMMENT-out.west);
\draw[-{Triangle[length=4bp,width=3bp]}] (STR_END) to[bend left=0] (STR_END-out.north);
\draw[-{Triangle[length=4bp,width=3bp]}] (WS) to[bend left=0] (WS-out.west);
\draw[-{Triangle[length=4bp,width=3bp]}] (START) |- node[pos=0.817,above,inner sep=1pt,align=center] {\code{LETTER} \code{EXP}} (IDENT);
\draw[-{Triangle[length=4bp,width=3bp]}] (START) |- node[pos=0.847,above,inner sep=1pt,align=center] {\code{DIGIT}} (INT);
\draw[-{Triangle[length=4bp,width=3bp]}] (START) to[bend left=0] node[pos=0.5,above,inner sep=1pt,align=center] {\code{EQ} \code{SIMPLE}} (OP_DONE);
\draw[-{Triangle[length=4bp,width=3bp]}] (START) |- node[pos=0.672,above,inner sep=1pt,align=center] {\code{MINUS}} (MINUS);
\draw[-{Triangle[length=4bp,width=3bp]}] (START) |- node[pos=0.661,above,inner sep=1pt,align=center] {\code{PLUS} \code{STAR} \code{LT} \code{GT}} (OP_EQ);
\draw[-{Triangle[length=4bp,width=3bp]}] (START) |- node[pos=0.756,above,inner sep=1pt,align=center] {\code{SLASH}} (SLASH);
\draw[-{Triangle[length=4bp,width=3bp]}] (START) |- node[pos=0.756,above,inner sep=1pt,align=center] {\code{DOT}} (DOT);
\draw[-{Triangle[length=4bp,width=3bp]}] (START) |- node[pos=0.854,above,inner sep=1pt,align=center] {\code{QUOTE}} (STR_BODY);
\draw[-{Triangle[length=4bp,width=3bp]}] (START) |- node[pos=0.878,above,inner sep=1pt,align=center] {\code{WS} \code{NL}} (WS);
\path[-{Triangle[length=4bp,width=3bp]}] (IDENT) edge[loop above,min distance=7mm] node[above,inner sep=1pt,align=center] {\code{LETTER} \code{EXP} \code{DIGIT} \code{\_}} (IDENT);
\path[-{Triangle[length=4bp,width=3bp]}] (INT) edge[loop above,min distance=7mm] node[above,inner sep=1pt,align=center] {\code{DIGIT}} (INT);
\draw[-{Triangle[length=4bp,width=3bp]}] (INT) to[bend left=0] node[pos=0.5,above,inner sep=1pt,align=center] {\code{DOT}} (INT_DOT);
\draw[-{Triangle[length=4bp,width=3bp]}] (INT) to[bend left=0] node[pos=0.5,below left,inner sep=1pt,align=center] {\code{EXP}} (EXP);
\draw[-{Triangle[length=4bp,width=3bp]}] (INT_DOT) to[bend left=0] node[pos=0.5,above,inner sep=1pt,align=center] {\code{DIGIT}} (FLOAT);
\path[-{Triangle[length=4bp,width=3bp]}] (FLOAT) edge[loop right,min distance=7mm] node[right,inner sep=1pt,align=center] {\code{DIGIT}} (FLOAT);
\draw[-{Triangle[length=4bp,width=3bp]}] (FLOAT) to[bend left=0] node[pos=0.5,right,inner sep=1pt,align=center] {\code{EXP}} (EXP);
\draw[-{Triangle[length=4bp,width=3bp]}] (EXP) to[bend left=0] node[pos=0.5,above left,inner sep=1pt,align=center] {\code{PLUS} \code{MINUS}} (EXP_SIGN);
\draw[-{Triangle[length=4bp,width=3bp]}] (EXP) to[bend left=0] node[pos=0.5,right,inner sep=1pt,align=center] {\code{DIGIT}} (FLOAT_EXP);
\draw[-{Triangle[length=4bp,width=3bp]}] (EXP_SIGN) to[bend left=0] node[pos=0.5,above,inner sep=1pt,align=center] {\code{DIGIT}} (FLOAT_EXP);
\path[-{Triangle[length=4bp,width=3bp]}] (FLOAT_EXP) edge[loop right,min distance=7mm] node[right,inner sep=1pt,align=center] {\code{DIGIT}} (FLOAT_EXP);
\draw[-{Triangle[length=4bp,width=3bp]}] (MINUS) -| node[pos=0.353,above,inner sep=1pt,align=center] {\code{GT} \code{EQ}} (OP_DONE);
\draw[-{Triangle[length=4bp,width=3bp]}] (OP_EQ) -| node[pos=0.364,above,inner sep=1pt,align=center] {\code{EQ}} (OP_DONE);
\draw[-{Triangle[length=4bp,width=3bp]}] (SLASH) -| node[pos=0.273,above,inner sep=1pt,align=center] {\code{EQ}} (OP_DONE);
\draw[-{Triangle[length=4bp,width=3bp]}] (SLASH) to[bend left=0] node[pos=0.5,left,inner sep=1pt,align=center] {\code{SLASH}} (COMMENT);
\path[-{Triangle[length=4bp,width=3bp]}] (COMMENT) edge[loop left,min distance=7mm] node[left,inner sep=1pt,align=center] {\code{any except NL}} (COMMENT);
\draw[-{Triangle[length=4bp,width=3bp]}] (DOT) -| node[pos=0.273,below,inner sep=1pt,align=center] {\code{STAR} \code{SLASH}} (OP_DONE);
\draw[-{Triangle[length=4bp,width=3bp]}] (DOT) to[bend left=0] node[pos=0.5,above left,inner sep=1pt,align=center] {\code{DOT}} (DOT_DOT);
\draw[-{Triangle[length=4bp,width=3bp]}] (DOT_DOT) -| node[pos=0.134,above,inner sep=1pt,align=center] {\code{DOT}} (OP_DONE);
\path[-{Triangle[length=4bp,width=3bp]}] (STR_BODY) edge[loop above,min distance=7mm] node[above,inner sep=1pt,align=center] {\code{any except NL QUOTE}} (STR_BODY);
\draw[-{Triangle[length=4bp,width=3bp]}] (STR_BODY) to[bend left=0] node[pos=0.5,above,inner sep=1pt,align=center] {\code{QUOTE}} (STR_END);
\path[-{Triangle[length=4bp,width=3bp]}] (WS) edge[loop above,min distance=7mm] node[above,inner sep=1pt,align=center] {\code{WS} \code{NL}} (WS);
\end{tikzpicture}
\end{adjustbox}
\caption{The merged DFA of the lexer.}
\label{fig:dfa-merged}
\end{figure}

\section{The Lexical Analyzer Algorithm}
\label{sec:lexer-algorithm}

The scanner exposes a single operation, \code{next}, which returns the next
token each time it is called. It keeps the source text, the current position,
and the current line and column, both counted from 1. The reference
implementation is a C++ program.

\subsection{Maximal Munch}

Each call to \code{next} starts at the current position and walks the
transition table for as long as a transition exists. Whenever it enters an
accepting state it records that state and the position reached. When it
reaches the dead state, or the end of the input, it rewinds to the last
recorded position. This is the longest match rule: the lexer always takes the
longest prefix that forms a token \cite{aho2006compilers}. The rewinding is
what allows one automaton to handle lexemes that share a prefix.

The text \code{0...2} shows this. The lexer reads \code{0}, which is accepting,
and then a dot, which moves to \code{INT\_DOT}. The next dot has no transition
from \code{INT\_DOT}, so the lexer rewinds to the \code{0} and emits it as an
integer literal. Scanning resumes at the first dot, which reaches
\code{OP\_DONE} after three dots and emits the ellipsis, and then \code{2} is
scanned on its own. In the same way, \code{1e+;} is scanned as the integer
\code{1}, the identifier \code{e}, and the operator \code{+}, because
\code{EXP} and \code{EXP\_SIGN} are not accepting.

\subsection{The Scanning Loop}

Listing~\ref{lst:scan-loop} shows the body of \code{next}. The character class
mapping is written \code{classOf}, and \code{advance} moves the position
forward while updating the line and column. The two lookup tables described
earlier are used through two functions. The function \code{lookupKeyword}
handles words, and \code{lookupOperator} handles operators.

\begin{lstlisting}[style=pseudocode, float=tbp,
    caption={The scanning loop of the lexer.}, label=lst:scan-loop]
function next():
    loop:
        if pos >= length(src):
            return Token(COMPILER_EOF)

        start <- pos
        state <- START
        lastAccept <- DEAD
        lastEnd <- start
        cur <- start

        while cur < length(src):
            nxt <- TRANSITION[state][classOf(src[cur])]
            if nxt is DEAD: break
            state <- nxt
            cur <- cur + 1
            if ACTION[state] is not NONE:
                lastAccept <- state
                lastEnd <- cur

        if lastAccept is DEAD:
            end <- start + 1 if cur = start else cur
            lexeme <- src[start .. end)
            advance(end - start)
            return Token(COMPILER_ERROR, lexeme)

        lexeme <- src[start .. lastEnd)
        advance(lastEnd - start)

        switch ACTION[lastAccept]:
            WORD:    type <- lookupKeyword(lexeme)
            INT:     type <- CONSTANTS_INT_LITERAL
            FLOAT:   type <- CONSTANTS_FLOAT_LITERAL
            OP:      type <- lookupOperator(lexeme)
            STRING:  type <- CONSTANTS_STRING_LITERAL
            SKIP:    continue    // scan again
        return Token(type, lexeme)
\end{lstlisting}

Each token also records the line and column where it started, which are taken
before the position moves.

\subsection{Resolving the Token}

The action of the last accepting state says how to name the lexeme. A word is
looked up in the keyword table, which also holds the type names, and falls back
to an identifier. An operator or punctuator is looked up in the operator
table. Integer literals, float literals, and string literals are named
directly by their state. Whitespace and comments are skipped, and the loop
starts again from the new position, so \code{next} never returns them.

The token set lists a bool constant token for the values \code{0} and
\code{1}. The lexer does not produce it. Both values are scanned as integer
literals, and interpreting them as booleans is left to the parser.

\subsection{Errors}

If the walk never reached an accepting state, the text it consumed becomes the
lexeme of an error token, and at least one character is consumed so that
scanning always makes progress. The lexer then reports the lexeme on the
standard error stream together with its line and column, and scanning
resumes right after it. One run therefore reports every lexical error in a
file. Table~\ref{tab:lexer-errors} shows how some invalid inputs are reported.

\begin{table}[tbp]
\centering
\small
\renewcommand{\arraystretch}{1.1}
\begin{tabular}{@{}ll@{}}
\toprule
\textbf{Input} & \textbf{Reported lexeme} \\
\midrule
\code{\_hidden} & \code{\_} \\
\code{1.} & \code{.} \\
\code{.5} & \code{.} \\
\code{1.5.2} & \code{.} \\
\code{3..4} & \code{..} \\
\code{"unterminated} & \code{"unterminated} \\
\code{!=} & \code{!} \\
\code{a @ b} & \code{@} \\
\bottomrule
\end{tabular}
\caption{Examples of lexical errors.}
\label{tab:lexer-errors}
\end{table}

\subsection{Token Output}

The lexer writes one token per line in the form \code{<TYPE, value>}, ending
with the end of file token. Only some tokens have a value part: identifiers,
scalar type names, integer, float, and string literals, the arithmetic
operators, and the assignment operators. These are the tokens whose Values
column is filled in Table~\ref{tab:token-set}. Every other token is written
with an empty value part, for example \code{<PUNCTUATION\_LBRACE,>}.

\subsection{Sample Runs}

To see the algorithm at work, we run the lexer on three statements. For each
lexeme, the tables list the states the automaton visits, starting from
\code{START}, followed by the lookup that names the lexeme when there is one,
and ending in the token that comes out. The whitespace between lexemes
takes the automaton through \code{WS}, which is skipped, so it does not appear
in the tables. After the last lexeme of the input, the lexer returns the end of
file token, \code{<COMPILER\_EOF,>}, which is also left out.

\medskip
\noindent\begin{minipage}{\textwidth}
\noindent\textbf{Statement 1.} The first statement declares a variable. Every lexeme is either a word, a number, or a single character.

\begin{lstlisting}
i32 count = 10;
\end{lstlisting}

\begin{center}
\footnotesize
\renewcommand{\arraystretch}{1.1}
\setlength{\tabcolsep}{4pt}
\begin{tabular}{@{}l>{\raggedright\arraybackslash}p{370pt}@{}}
\toprule
\textbf{Lexeme} & \textbf{States visited, lookup, and token} \\
\midrule
\code{i32} & \tok{START} $\to$ \tok{IDENT} $\to$ \tok{IDENT} $\to$ \tok{IDENT} $\to$ lookup(keyword table) $\to$ \tok{<DT_SCALAR, i32>} \\
\code{count} & \tok{START} $\to$ \tok{IDENT} $\to$ \tok{IDENT} $\to$ \tok{IDENT} $\to$ \tok{IDENT} $\to$ \tok{IDENT} $\to$ lookup(keyword table) $\to$ \tok{<IDENTIFIER, count>} \\
\code{=} & \tok{START} $\to$ \tok{OP_DONE} $\to$ lookup(operator table) $\to$ \tok{<OPERATORS_ASSIGNMENT, =>} \\
\code{10} & \tok{START} $\to$ \tok{INT} $\to$ \tok{INT} $\to$ \tok{<CONSTANTS_INT_LITERAL, 10>} \\
\code{;} & \tok{START} $\to$ \tok{OP_DONE} $\to$ lookup(operator table) $\to$ \tok{<PUNCTUATION_SEMICOLON,>} \\
\bottomrule
\end{tabular}
\end{center}
\end{minipage}
\medskip


\medskip
\noindent\begin{minipage}{\textwidth}
\noindent\textbf{Statement 2.} The second statement has a float with an exponent, and a matrix operator that is written without spaces around it.

\begin{lstlisting}
f32 y = 2.5e-3 + a.*b;
\end{lstlisting}

\begin{center}
\footnotesize
\renewcommand{\arraystretch}{1.1}
\setlength{\tabcolsep}{4pt}
\begin{tabular}{@{}l>{\raggedright\arraybackslash}p{370pt}@{}}
\toprule
\textbf{Lexeme} & \textbf{States visited, lookup, and token} \\
\midrule
\code{f32} & \tok{START} $\to$ \tok{IDENT} $\to$ \tok{IDENT} $\to$ \tok{IDENT} $\to$ lookup(keyword table) $\to$ \tok{<DT_SCALAR, f32>} \\
\code{y} & \tok{START} $\to$ \tok{IDENT} $\to$ lookup(keyword table) $\to$ \tok{<IDENTIFIER, y>} \\
\code{=} & \tok{START} $\to$ \tok{OP_DONE} $\to$ lookup(operator table) $\to$ \tok{<OPERATORS_ASSIGNMENT, =>} \\
\code{2.5e-3} & \tok{START} $\to$ \tok{INT} $\to$ \tok{INT_DOT} $\to$ \tok{FLOAT} $\to$ \tok{EXP} $\to$ \tok{EXP_SIGN} $\to$ \tok{FLOAT_EXP} $\to$ \tok{<CONSTANTS_FLOAT_LITERAL, 2.5e-3>} \\
\code{+} & \tok{START} $\to$ \tok{OP_EQ} $\to$ lookup(operator table) $\to$ \tok{<OPERATORS_STANDARD_PLUSMINUS, +>} \\
\code{a} & \tok{START} $\to$ \tok{IDENT} $\to$ lookup(keyword table) $\to$ \tok{<IDENTIFIER, a>} \\
\code{.*} & \tok{START} $\to$ \tok{DOT} $\to$ \tok{OP_DONE} $\to$ lookup(operator table) $\to$ \tok{<OPERATORS_MATRIX_DOT_STAR,>} \\
\code{b} & \tok{START} $\to$ \tok{IDENT} $\to$ lookup(keyword table) $\to$ \tok{<IDENTIFIER, b>} \\
\code{;} & \tok{START} $\to$ \tok{OP_DONE} $\to$ lookup(operator table) $\to$ \tok{<PUNCTUATION_SEMICOLON,>} \\
\bottomrule
\end{tabular}
\end{center}
\end{minipage}
\medskip


\medskip
\noindent\begin{minipage}{\textwidth}
\noindent\textbf{Statement 3.} The third statement is the range loop, and the range \code{1...5} is the case where the lexer has to rewind.

\begin{lstlisting}
for (i32 i = 1...5) {
\end{lstlisting}

\begin{center}
\footnotesize
\renewcommand{\arraystretch}{1.1}
\setlength{\tabcolsep}{4pt}
\begin{tabular}{@{}l>{\raggedright\arraybackslash}p{370pt}@{}}
\toprule
\textbf{Lexeme} & \textbf{States visited, lookup, and token} \\
\midrule
\code{for} & \tok{START} $\to$ \tok{IDENT} $\to$ \tok{IDENT} $\to$ \tok{IDENT} $\to$ lookup(keyword table) $\to$ \tok{<KW_LOOPS_FOR,>} \\
\code{(} & \tok{START} $\to$ \tok{OP_DONE} $\to$ lookup(operator table) $\to$ \tok{<PUNCTUATION_LPAREN,>} \\
\code{i32} & \tok{START} $\to$ \tok{IDENT} $\to$ \tok{IDENT} $\to$ \tok{IDENT} $\to$ lookup(keyword table) $\to$ \tok{<DT_SCALAR, i32>} \\
\code{i} & \tok{START} $\to$ \tok{IDENT} $\to$ lookup(keyword table) $\to$ \tok{<IDENTIFIER, i>} \\
\code{=} & \tok{START} $\to$ \tok{OP_DONE} $\to$ lookup(operator table) $\to$ \tok{<OPERATORS_ASSIGNMENT, =>} \\
\code{1} & \tok{START} $\to$ \tok{INT} $\to$ \tok{INT_DOT} \textit{(no move on the next dot, rewind to \tok{INT})} $\to$ \tok{<CONSTANTS_INT_LITERAL, 1>} \\
\code{...} & \tok{START} $\to$ \tok{DOT} $\to$ \tok{DOT_DOT} $\to$ \tok{OP_DONE} $\to$ lookup(operator table) $\to$ \tok{<PUNCTUATION_ELLIPSIS,>} \\
\code{5} & \tok{START} $\to$ \tok{INT} $\to$ \tok{<CONSTANTS_INT_LITERAL, 5>} \\
\code{)} & \tok{START} $\to$ \tok{OP_DONE} $\to$ lookup(operator table) $\to$ \tok{<PUNCTUATION_RPAREN,>} \\
\code{\{} & \tok{START} $\to$ \tok{OP_DONE} $\to$ lookup(operator table) $\to$ \tok{<PUNCTUATION_LBRACE,>} \\
\bottomrule
\end{tabular}
\end{center}
\end{minipage}
\medskip


\medskip
\noindent In the first two runs, every lexeme is accepted at the last state it reaches,
and the lexer never has to go back. The word \code{i32} ends in \code{IDENT},
and the keyword table turns it into a type token. The float \code{2.5e-3} is
the longest path in the automaton, and it passes through three states that
are not accepting. In the third run, the \code{1} is followed by a dot, which
moves the automaton to \code{INT\_DOT}. The next character is another dot, which
has no transition there, so the lexer goes back to the last accepting state,
which was reached after the \code{1}, and returns the integer. The scan for the
next token then starts at the first dot, and reads all three of them.

\section{Lexers in the Real World}
\label{sec:lexers-real-world}

The lexer of OctaC is deliberately small. This section looks at how lexers
are built in practice, what they handle that ours does not, and which
simplifications we made.

\subsection{Lexer Generators}

Many lexers are produced by a generator such as Lex \cite{lesk1975lex} or its
successor Flex \cite{flex-manual}. The programmer writes a list of regular
expressions, each with an action. The generator converts them into a
deterministic finite automaton and emits a table-driven scanner
\cite{aho2006compilers}. Two rules resolve conflicts between patterns: the
longest match wins, and among patterns that match the same length, the one
listed first wins \cite{lesk1975lex,flex-manual}. Our scanner follows the
longest match rule. The second rule is not needed, because keywords are not
patterns of their own.

Flex recommends matching identifiers with one pattern and looking keywords up
in a table, instead of writing a pattern for each keyword
\cite{flex-manual}, which is what OctaC does. Flex also warns that a scanner
which has to back up over several characters is slower, and offers ways to
find and remove such cases \cite{flex-manual}. Our scanner backs up too, as the
run of \code{1...5} shows, but only over a few characters. Generators also
compress their tables \cite{aho2006compilers}. Ours has 19 states and 17
character classes, which is small enough to be stored as it is.

\subsection{Hand-Written Lexers}

Many production compilers do not use a generator. The lexer of the GNU C
preprocessor is coded by hand and is not implemented as a state machine
\cite{gcc-cppinternals}, and Clang has a hand-written lexer that scans the
source buffer directly \cite{clang-internals}. A hand-written lexer gives full
control over the speed of the scanner, the quality of error messages, and the
handling of features that regular expressions describe poorly. The lexer of
OctaC sits between the two approaches. The table is written by hand, but the
scanning loop is the same one that a generator would emit.

\subsection{What Real Lexers Handle}

Lexers for established languages deal with many things that OctaC leaves out:

\begin{itemize}
    \item \textbf{Richer literals.} Numbers may be written in hexadecimal,
    octal, or binary, with digit separators and type suffixes. Strings have
    escape sequences, raw forms, and may span lines \cite{rust-reference}.
    \item \textbf{More kinds of comments.} Block comments exist in many
    languages, and in Rust they can be nested \cite{rust-reference}.
    \item \textbf{Unicode.} Identifiers may contain letters from any script,
    following the rules of Unicode Standard Annex 31
    \cite{unicode-uax31,python-reference}.
    \item \textbf{Layout.} Python's lexer produces indent and dedent tokens
    from the indentation of each line \cite{python-reference}. The lexer of
    Go inserts semicolons at the ends of lines by a fixed rule, so programs
    can leave most of them out \cite{go-spec}.
    \item \textbf{Preprocessing.} A C lexer works together with the
    preprocessor, which handles macros and file inclusion
    \cite{gcc-cppinternals}.
\end{itemize}

\subsection{Simplifications in OctaC}

Table~\ref{tab:lexer-simplifications} summarizes the differences. Each of the
simplifications keeps the automaton small, and each could be lifted later by
adding states and character classes to the table.

\begin{table}[tbp]
\centering
\small
\renewcommand{\arraystretch}{1.15}
\begin{tabular}{@{}>{\raggedright\arraybackslash}p{62pt}>{\raggedright\arraybackslash}p{190pt}>{\raggedright\arraybackslash}p{180pt}@{}}
\toprule
\textbf{Concern} & \textbf{Real-world lexers} & \textbf{OctaC} \\
\midrule
Construction & Generated from regular expressions, or written by hand \cite{flex-manual,clang-internals} & One hand-written transition table \\
Characters & Unicode text, usually as UTF-8 & Bytes; anything outside the listed classes is an error \\
Identifiers & Unicode letters allowed \cite{unicode-uax31} & ASCII letters, digits, and underscores, starting with a letter \\
Numbers & Hexadecimal, binary, separators, suffixes \cite{rust-reference} & Decimal integers and floats, with an optional exponent \\
Strings & Escapes, raw strings, multiple lines \cite{rust-reference} & None of these: one line, no escapes \\
Comments & Block comments, sometimes nested \cite{rust-reference} & Line comments with \code{//} only \\
Layout & Indentation tokens, semicolon insertion \cite{python-reference,go-spec} & Whitespace is not significant, and semicolons are explicit \\
Preprocessing & Macros and file inclusion \cite{gcc-cppinternals} & None \\
Keywords & Table lookup after matching an identifier \cite{flex-manual} & Same \\
Errors & Diagnostics with recovery & Error token, a diagnostic, and scanning resumes \\
Input & Streamed through a buffer & The whole file is read into memory \\
\bottomrule
\end{tabular}
\caption{Real-world lexers compared with the OctaC lexer.}
\label{tab:lexer-simplifications}
\end{table}
